\documentclass[10pt,a4paper]{amsart}
\usepackage[margin=19mm]{geometry}
\usepackage{amsmath,amssymb,mathtools}
\usepackage[scr=rsfs,cal=euler]{mathalfa}
\usepackage{xfrac}
\usepackage[unicode,hidelinks,bookmarks=false]{hyperref}
\newcommand{\ie}{i.e.\ }
\newcommand{\cf}{cf.\ }
\newcommand{\resp}{respectively\ }
\newcommand{\Subsection}{\subsection}
\providecommand{\nobreakdash}{\nobreak\hbox{-}}
\setlength{\emergencystretch}{2em}
\multlinegap0pt
\usepackage{amssymb}
\newtheorem{thm}{Theorem}
\newtheorem{pf}{Pf}
\newcommand{\norm}[1]{\left\lVert#1\right\rVert}
\title{Arithmetic Bohr radius for the Minkowski space}
\author{Vasudevarao Allu, Himadri Halder,, Subhadip Pal}
\date{Source: arXiv:2309.15550v3 (CC BY 4.0)}
\begin{document}
\maketitle

\noindent\emph{Source and licence.} Arithmetic Bohr radius for the Minkowski space. Version arXiv:2309.15550v3 (CC BY 4.0), distributed by arXiv under the Creative Commons Attribution 4.0 International licence, http://creativecommons.org/licenses/by/4.0/, as recorded in the arXiv licence metadata for this exact version and shown on the abstract page https://arxiv.org/abs/2309.15550v3. Full licence notice: \texttt{licenses/arxiv-2309.15550-CC-BY-4.0.txt}.\par\smallskip

\noindent\textit{Editorial scope.} Counted unit: the article's thm Pal-Vasu-P4-thm-2.3 (source0.tex lines 373--378). The article's complete original proof of it is delivered with it (source0.tex lines 486--527, one \texttt{proof} environment of 42 lines, which the article places away from the statement). No mathematics is written, repaired or re-derived: the statement and the proof are the article's own lines, in the article's own notation, and no retained line is substituted or re-flowed.  No further source-local context is needed: every cross-reference of every retained line resolves inside the two delivered spans.  Nothing was omitted from any retained span, so the package carries no omission record.  No retained line contains a citation, so the wrapper carries no bibliography and no bibliography entry.  No retained line contains a figure environment, an image-inclusion command or drawing code, so no image is delivered and the package declares no asset dependency.  Editorial formatting only, in addition: an AMS \texttt{amsart} preamble that restates the article's own declarations and macros used by the retained lines, namely the environments \texttt{thm}, \texttt{pf} on separate counters, together with the standard packages those lines need.  The retained environments print in this document as Theorem 1, Pf 1 in order of appearance, whereas the article prints the same items with the numbering of its own full document.  The complete native source of the article as distributed in the arXiv e-print for this exact version is bundled unchanged as \texttt{sources/147632/source0.tex}.
\begin{thm}\label{Pal-Vasu-P4-thm-2.3}
Let $1\leq p<\infty$. Then for every $n\in \mathbb{N}$ and $1\leq q<\infty$ we have 
\begin{equation*}
	\frac{H^1_p(\mathbb{D})}{n}\leq A_p\left(B_{\ell^n_q}\right)\leq \left(\frac{H^1_p(\mathbb{D})}{n^{1/p}}\right)^{1/q}.
\end{equation*}
\end{thm}

\begin{pf} [{\bf Proof of Theorem \ref{Pal-Vasu-P4-thm-2.3}}]
	First we show the left-hand inequality
	\begin{equation*}
		\frac{H^1_p(\mathbb{D})}{n}\leq A_p\left(B_{\ell^n_q}\right).
	\end{equation*}
	Assume $r=H^1_p(\mathbb{D})$ and $f\in \mathcal{B}(B_{\ell^n_q})$. We define $g(z)=f(ze_1)=f(z,0,.\,.\,.,0)$ for $z\in \mathbb{D}.$ Then $g:\mathbb{D}\rightarrow\mathbb{C}$ will be a holomorphic function on $\mathbb{D}$ with $\mbox{Re}(g(z))>0$ and $g(0)=1$. Therefore,
	\begin{align*}
		\frac{1}{2}\left\{|c_0(f)|^p+\sum_{k=1}^{\infty}\sum_{|\alpha|=k}|c_{\alpha}(f)|(r,0,.\,.\,.\,,0)^{p\alpha}\right\}^{\frac{1}{p}}= \frac{1}{2}\left\{|c_0(g)|^p+\sum_{k=1}^{\infty}|c_{k}(g)|r^{pk}\right\}^{\frac{1}{p}}\leq 1
	\end{align*}
	for all $f(z)=\sum_{\alpha\in \mathbb{N}^n_{0}}c_{\alpha}(f)z^{\alpha} \in \mathcal{B}(B_{\ell^n_q})$. Hence, we obtain $r/n\leq A_p(B_{\ell^n_q})$, which gives our desired inequality.\\[2mm]
	On the other hand, we want to prove that 
	\begin{equation*}
		A_p\left(B_{\ell^n_q}\right)\leq \left(\frac{H^1_p(\mathbb{D})}{n^{1/p}}\right)^{1/q}.
	\end{equation*}
	Let $r\in \mathbb{R}^n_{\geq 0}$ be such that for all $u\in \mathcal{B}(B_{\ell^n_q})$, we have
	\begin{equation*}
		\frac{1}{2}\left\{|c_0(u)|^p + \sum_{k=1}^{\infty}\sum_{|\alpha|=k}|c_{\alpha}(u)|^pr^{p\alpha}\right\}^{\frac{1}{p}}\leq 1.
	\end{equation*}
	Now it is sufficient to show that
	\begin{equation}\label{Pal-Vasu-P4-e-3.5}
		\frac{1}{n}\left(\sum_{j=1}^{n}r_j\right) \leq  \left(\frac{H^1_p(\mathbb{D})}{n^{1/p}}\right)^{1/q}.
	\end{equation}
	Fix $f\in \mathcal{B}(\mathbb{D})$, and define the function
	\begin{equation*}
		v(z)=z^q_1+\cdots+z^q_n, \quad z=(z_1,.\,.\,.,z_n)\in B_{\ell^n_q}. 
	\end{equation*}
	 Consider $u=f\circ v$ such that $\mbox{Re}(u(z))=\mbox{Re}(f(v(z)))>0$ and $u(0)=f(v(0))=f(0)=1$. Moreover, for each $z\in B_{\ell^n_q}$, we have 
	\begin{equation*}
		u(z)=\sum_{k=0}^{\infty}c_{k}(f)v(z)^k=\sum_{k=0}^{\infty}c_k(f)\sum_{|\alpha|=k}\frac{k!}{\alpha!}z^{q\alpha}=\sum_{\alpha\in \mathbb{N}^n_{0}}c_{\alpha}(u)z^{q\alpha},
	\end{equation*}
	where $c_{\alpha}(u)=c_k(f)(k!/\alpha!)$ whenever $|\alpha|=k$. Then for all $z\in B_{\ell^n_q}$, we have
	\begin{align*}
		\frac{1}{2}\left\{|c_0(u)|^p + \sum_{k=1}^{\infty}\sum_{|\alpha|=k}|c_{\alpha}(u)z^{q\alpha}|^p\right\}^{\frac{1}{p}} &= \frac{1}{2}\left\{|c_0(f)|^p+\sum_{k=1}^{\infty}|c_{k}(f)|^p\sum_{|\alpha|=k}\left(\frac{k!}{\alpha!}\right)^p|z|^{pq\alpha}\right\}^{\frac{1}{p}}\\& \geq 
		\frac{1}{2}\left\{|c_0(f)|^p+\sum_{k=1}^{\infty}|c_{k}(f)|^p\sum_{|\alpha|=k}\frac{k!}{\alpha!}|z|^{pq\alpha}\right\}^{\frac{1}{p}}\\ & = 
		\frac{1}{2}\left\{|c_0(f)|^p+\sum_{k=1}^{\infty}|c_{k}(f)|^p\norm{z}^{pqk}_{pq}\right\}^{\frac{1}{p}}
	\end{align*}
	so that finally we have 
	\begin{align*}
		\frac{1}{2}\left\{|c_0(f)|^p+\sum_{k=1}^{\infty}|c_{k}(f)|^p\norm{r}^{pqk}_{pq}\right\}^{\frac{1}{p}}\leq \frac{1}{2}\left\{|c_0(u)|^p + \sum_{k=1}^{\infty}\sum_{|\alpha|=k}|c_{\alpha}(u)r^{q\alpha}|^p\right\}^{\frac{1}{p}}\leq 1. 
	\end{align*}
	It follows that $\norm{r}^q_{pq}\leq H^1_{p}(\mathbb{D})$. By virtue of H\"{o}lder's inequality, we have $\norm{r}^{pq}_{1}\leq n^{pq-1}\norm{r}^{pq}_{pq}$. Hence, we obtain $n^{1-pq}\norm{r}^{pq}_{1}\leq (H^1_{p}(\mathbb{D}))^p$, which gives the estimate \eqref{Pal-Vasu-P4-e-3.5}. This completes the proof. 
\end{pf}

\end{document}
